\chapter{Reproducibility}
\label{app:repro}

\section{Computing environment}
The work began on Windows 11 using the Windows Subsystem for Linux (WSL2). WSL2
by default makes only about half of the system memory available to Linux,
roughly 9\,GB of the 19\,GB here, which is less than the 9.3\,GB LOFAR file. On
my supervisor's advice the project was moved to a native Linux installation.
(Raising WSL2's memory limit in its configuration file would also have solved
this particular problem.) \Cref{tab:env} lists the final environment.

\begin{table}[H]
  \centering
  \caption{Computing environment.}
  \label{tab:env}
  \begin{tabular}{ll}
    \toprule
    Component & Version \\
    \midrule
    Operating system & Fedora 44, kernel 6.19.10 \\
    GPU & NVIDIA RTX 3060 Laptop, 6\,GB \\
    Driver / CUDA & 610.57.04 / 13.3 \\
    \tfunet{} environment & Python 3.12, TensorFlow 2.21.0 \\
    \hybrid{} environment & Python 3.14, PyTorch 2.13.0 \\
    \bottomrule
  \end{tabular}
\end{table}

TensorFlow and PyTorch need incompatible versions of NVIDIA's cuDNN library, so
they are kept in two separate Python environments. Training speed depends
strongly on power: on battery the GPU is throttled to 210\,MHz of its 2100\,MHz,
and one \tfunet{} training step took 2684\,ms instead of 174\,ms---15 times
slower.

\section{Code and data}
All code, analysis scripts and per-run metrics are in the project repository,
\url{https://github.com/RonakSinghRaina/MP}. The HERA and LOFAR datasets are
available from Zenodo\footcite{mesarcik2022data}; the synthetic datasets are
regenerated from the generator with seed 42. The main experiments are run with:
\begin{lstlisting}[language=bash]
# tf_unet on LOFAR (TensorFlow environment)
~/tf-env/bin/python experiments/lofar_tfunet_baseline.py --lr 1e-4 --epochs 150 --seed 0

# HybridRFINet on LOFAR (PyTorch environment)
~/torch-env/bin/python experiments/lofar_hybrid.py --base 8 --no_class_weight --seed 0

# Ablation: any variant through the same pipeline
~/torch-env/bin/python experiments/lofar_hybrid.py --variant plain_unet \
    --base 8 --no_class_weight --seed 0
\end{lstlisting}
